\documentclass{article}
\usepackage{amsmath,amsthm,amssymb}
\usepackage[margin=1in]{geometry}

\DeclareMathOperator{\area}{area}
\DeclareMathOperator{\dinv}{dinv}
\DeclareMathOperator{\defc}{defc}

\newtheorem{proposition}{Proposition}
\newtheorem{theorem}{Theorem}
\newtheorem{conjecture}{Conjecture}

\title{Flat Middle Coefficients for $q,t$-Catalan Polynomials}
\author{}
\date{}

\begin{document}
\maketitle

\section{Overview}

This note records a structural observation about the coefficients of
\(q,t\)-Catalan polynomials and their rational analogues. As in the discussion
of string decompositions, we consider three layers of generality: the classical
Dyck case, the family \(r=\tau s+1\), and the general rational case. In the
classical case, the observation follows from the proved low-deficit
skeleton-string decomposition. In the \(r=\tau s+1\) family, it follows
wherever the special-skeleton formula is accepted and the relevant skeleton
area bound is verified. In the general rational case, we treat the phenomenon
computationally by directly enumerating rational Dyck paths and checking the
coefficient slices.

The phenomenon is the following. Fix a total degree \(M\).  In a fixed deficit
layer \(d\), the monomials have the form
\[
  q^j t^{M-d-j}.
\]
The middle interval of this layer is
\[
  d\le j\le M-2d.
\]
The flat-middle assertion says that the coefficient of
\(q^j t^{M-d-j}\) is independent of \(j\) throughout this interval.  The
interval is nonempty exactly when \(3d\le M\).  When it is nonempty, it is the
part of the deficit layer left after discarding the first \(d\) possible
area-slices and the last \(d\) possible area-slices.

\section{Why string decompositions imply flat middle coefficients}

The useful mechanism is elementary.  Suppose a deficit-\(d\) layer is described
by strings.  A root object \(R\) of area \(a\) contributes the interval
\[
  q^a t^{M-d-a},\quad
  q^{a+1}t^{M-d-a-1},\quad\ldots,\quad
  q^{M-d-a}t^a.
\]
Equivalently, it contributes one copy of every monomial
\[
  q^j t^{M-d-j}
  \qquad(a\le j\le M-d-a).
\]
Thus a single root contributes to every middle slice
\(d\le j\le M-2d\) as soon as
\[
  a\le d.
\]
This gives the basic criterion.

\begin{proposition}[Root-area criterion]
Let a deficit-\(d\) layer of total degree \(M-d\) be expressed as a disjoint
sum of interval strings indexed by roots \(R\).  If every root satisfies
\[
  \area(R)\le d,
\]
then the coefficients of \(q^j t^{M-d-j}\) are constant for
\[
  d\le j\le M-2d.
\]
The common value is the number of roots in deficit \(d\).
\end{proposition}

\begin{proof}
A root \(R\) of area \(a\) contributes to the slice \(j\) exactly when
\[
  a\le j\le M-d-a.
\]
If \(a\le d\) and \(d\le j\le M-2d\), then the left inequality is immediate,
and the right inequality follows from
\[
  j\le M-2d\le M-d-a.
\]
So every root contributes exactly once to every middle slice.  Since the
strings are disjoint, the common coefficient is the number of roots.
\end{proof}

The rest of the note explains how this criterion is used in the three settings
considered here.

\section{Classical Dyck paths}

An ordinary Dyck sequence of length \(n\) is a sequence
\[
  x=(x_0,\ldots,x_{n-1})
\]
with
\[
  x_0=0,\qquad x_i\ge0,\qquad x_{i+1}\le x_i+1.
\]
Its area and dinv statistics are
\[
  \area(x)=\sum_i x_i
\]
and
\[
  \dinv(x)=
  \#\{(i,j):0\le i<j<n,
      \ x_i=x_j\text{ or }x_i=x_j+1\}.
\]
The total degree is
\[
  M=\binom n2,
\]
and the deficit is
\[
  \defc(x)=M-\area(x)-\dinv(x).
\]
The classical \(q,t\)-Catalan polynomial in this convention is
\[
  C_n(q,t)=\sum_x q^{\area(x)}t^{\dinv(x)},
\]
where the sum ranges over all ordinary Dyck sequences of length \(n\).

The 2026 paper proves a skeleton-string decomposition in the low-deficit range
\[
  \defc\le 2n-8.
\]
The roots of these strings are special Dyck skeletons.  The same paper proves
an area bound for full Dyck skeletons:
\[
  \area(S)\le \defc(S).
\]
Since special Dyck skeletons are full Dyck skeletons, this applies to every
classical string root in the decomposition.

\begin{theorem}[Classical flat middle coefficients]
For \(n\ge4\), let \(M=\binom n2\).  If
\[
  0\le d\le 2n-8,
\]
then the coefficient of
\[
  q^j t^{M-d-j}
\]
in \(C_n(q,t)\) is independent of \(j\) for
\[
  d\le j\le M-2d.
\]
The common value is the number of special Dyck skeletons of length \(n\) and
classical deficit \(d\).
\end{theorem}

This is not merely computational evidence.  It is a direct consequence of the
proved string decomposition and the proved skeleton area bound.  The code
included with this item only supplies finite regression checks of this proved
phenomenon.

\section{The family $r=\tau s+1$}

Now fix \(s\ge1\) and \(\tau>1\).  A normalized \(\tau\)-Dyck sequence of
length \(s\) is a sequence
\[
  x=(x_0,\ldots,x_{s-1})
\]
with
\[
  x_0=0,\qquad x_i\ge0,\qquad x_{i+1}\le x_i+\tau.
\]
The area statistic is again
\[
  \area(x)=\sum_i x_i.
\]
For a finite sequence \(x\), define
\[
  \dinv_\tau(x)=\sum_{i<j} d_\tau(x_i,x_j),
\]
where
\[
  d_\tau(a,b)=
  \begin{cases}
    \max(0,a+\tau-b),& a\le b,\\
    \max(0,b+1+\tau-a),& a>b.
  \end{cases}
\]
The total degree is
\[
  M_{s,\tau}=\tau\binom{s}{2},
\]
and
\[
  \defc_\tau(x)=M_{s,\tau}-\area(x)-\dinv_\tau(x).
\]

The preceding skeleton-string item formulates a special-skeleton interval
formula for this \(r=\tau s+1\) family.  The conjectural low-deficit range is
\[
  B=(s-2)(\tau+1)-4.
\]
If the special-skeleton formula holds through deficit \(B\), then the
flat-middle result follows from the same root-area criterion, provided every
special rational skeleton \(S\) in the checked range satisfies
\[
  \area(S)\le \defc_\tau(S).
\]

Thus this item separates two issues.  The previous item tests or records the
special-skeleton formula itself.  The present item checks the additional
condition needed to extract the flat-middle consequence from that formula.

\begin{proposition}[Conditional \(r=\tau s+1\) consequence]
Assume the special-skeleton interval formula for normalized \(\tau\)-Dyck
sequences of length \(s\) is valid through deficit \(B\).  If every special
rational skeleton \(S\) with \(\defc_\tau(S)\le B\) satisfies
\[
  \area(S)\le \defc_\tau(S),
\]
then for every \(0\le d\le B\), the coefficient of
\[
  q^j t^{M_{s,\tau}-d-j}
\]
is independent of \(j\) throughout
\[
  d\le j\le M_{s,\tau}-2d.
\]
The common value is the number of special rational skeletons of length \(s\)
and deficit \(d\).
\end{proposition}

The accompanying code is designed to match the finite grid used in the
skeleton-string item.  In particular, its official \(r=\tau s+1\) mode checks
the skeleton area inequality on the same parameter grid:
\[
\begin{array}{c|c}
\tau & \text{lengths }s\\
\hline
2 & 1\le s\le 14\\
3 & 1\le s\le 12\\
4 & 1\le s\le 10\\
5 & 1\le s\le 9.
\end{array}
\]
It does not re-check the special-skeleton formula; it checks only the extra
root-area inequality needed for flat middle coefficients.

\section{General rational slopes}

For a general coprime rational slope \(r/s\), we do not currently have a
skeleton-string decomposition of the same kind.  Therefore the code checks the
flat-middle phenomenon directly.

The model uses position coordinates.  Put
\[
  H_i=\left\lfloor \frac{ri}{s}\right\rfloor
  \qquad(0\le i<s),
\]
and let
\[
  M=\sum_{i=0}^{s-1}H_i.
\]
A rational Dyck path is represented by a sequence
\[
  Q=(Q_0,\ldots,Q_{s-1})
\]
with \(Q_0=0\), \(0\le Q_i\le H_i\), and path heights
\[
  P_i=H_i-Q_i
\]
weakly increasing.  The code uses
\[
  \area(Q)=\sum_i Q_i
\]
and the same rational deficit statistic used in the rational-path diagnostics.
Then
\[
  \dinv(Q)=M-\area(Q)-\defc(Q).
\]
For each checked slope, the code enumerates all valid paths, groups them by
\((\area,\dinv)\), and verifies directly that for each nonempty middle band
\[
  0\le d\le \left\lfloor\frac{M}{3}\right\rfloor,
\]
the coefficients of
\[
  q^j t^{M-d-j}
  \qquad(d\le j\le M-2d)
\]
are all equal.

This is a direct coefficient check, not a string-decomposition proof.  It is
included because it gives evidence that the flat-middle phenomenon extends
beyond the cases where roots and strings are currently available.

\section{What the code checks}

The companion file \texttt{code.py} has three independent layers.

\subsection*{Classical checks}

The classical check enumerates ordinary Dyck sequences for a bounded range of
\(n\).  It verifies three things in the proved low-deficit range
\(0\le d\le2n-8\):
\begin{itemize}
\item the direct coefficients are flat on \(d\le j\le M-2d\);
\item the common coefficient equals the number of special Dyck skeletons of
      deficit \(d\);
\item every special skeleton encountered satisfies \(\area(S)\le\defc(S)\).
\end{itemize}
The default run checks \(4\le n\le8\).  This is a regression check of a proved
classical theorem, not the source of the theorem.

\subsection*{The $r=\tau s+1$ area-bound checks}

The \(r=\tau s+1\) check enumerates normalized \(\tau\)-Dyck sequences in the
retained deficit range and filters for special rational skeletons.  For each
such skeleton it verifies
\[
  \area(S)\le \defc_\tau(S).
\]
This is the extra condition needed to turn a special-skeleton formula into a
flat-middle coefficient statement.  The command
\[
  \texttt{python code.py --r1mod-mode official}
\]
runs the area-bound check on the official grid displayed above.  The default
mode runs a smaller representative subset so that the file gives a quick smoke
test.

\subsection*{Direct rational checks}

The direct rational check does not use skeletons.  It enumerates all rational
Dyck paths for the requested slope \(r/s\), forms the coefficient dictionary,
and checks every nonempty middle band.  The default slopes are
\[
  \frac53,
  \frac74,
  \frac85,
  \frac{10}{7},
  \frac{11}{8},
  \frac{13}{8}.
\]
A broader grid run used for this note checked every coprime slope
\[
  \frac{r}{s}
  \qquad\text{with}\qquad
  2\le s\le8,
  \quad s<r\le13,
\]
through all nonempty middle bands \(0\le d\le\lfloor M/3\rfloor\).  The
corresponding command is
\begin{verbatim}
python code.py --skip-classical --r1mod-mode none --rational-grid \
  --rational-s-max 8 --rational-r-max 13
\end{verbatim}
All of these direct rational checks are finite computational evidence.
They should not be read as a proof of the general rational conjecture.

\section{Status summary}

\begin{center}
\begin{tabular}{p{0.28\linewidth}|p{0.58\linewidth}}
case & status\\
\hline
classical Dyck & proved for \(0\le\defc\le2n-8\), using the 2026
skeleton-string decomposition and the skeleton area bound\\
\(r=\tau s+1\) & follows in any checked or accepted special-skeleton-formula
range once the special skeletons satisfy \(\area\le\defc\); the accompanying
code checks this root-area condition on the matching finite grid\\
general rational & checked directly by finite enumeration for selected slopes;
no general skeleton decomposition is claimed here
\end{tabular}
\end{center}

\begin{thebibliography}{9}

\bibitem{Hawkes2026DyckSymmetric}
Graham Hawkes,
\emph{Dyck Symmetric Functions and Applications to \(q,t\)-Catalan
Polynomials}, arXiv:2605.13003, 2026.

\end{thebibliography}

\end{document}
